\documentclass[11pt,a4paper]{article}
\pdfvariable trailerid {[
  <2BA4891E65C7A4C827CFA11C906CE6B0>
  <2BA4891E65C7A4C827CFA11C906CE6B0>
]}
\pdfextension catalog {/Lang (en-US)}

\usepackage[english]{babel}
\usepackage{amsmath,amssymb,amsthm,mathtools}
\usepackage{booktabs,enumitem,microtype,tabularx}
\usepackage[a4paper,margin=28mm]{geometry}
\usepackage{xcolor}
\usepackage{xurl}
\usepackage[colorlinks=true,linkcolor=blue!55!black,citecolor=blue!55!black,
 urlcolor=blue!55!black]{hyperref}
\Urlmuskip=0mu plus 2mu\relax

\hypersetup{
  pdftitle={Symmetry premiums, cyclic-closure coverage, and defect-spectrum persistence in Kautz digraphs},
  pdfauthor={Oleksiy Babanskyy},
  pdfsubject={Invariant feedback vertex sets under reverse-complement symmetries},
  pdfkeywords={Kautz digraphs, feedback vertex sets, involutions, cyclic closure, defect spectrum}
}

\newtheorem{theorem}{Theorem}[section]
\newtheorem{proposition}[theorem]{Proposition}
\newtheorem{lemma}[theorem]{Lemma}
\newtheorem{corollary}[theorem]{Corollary}
\theoremstyle{remark}
\newtheorem{remark}[theorem]{Remark}

\newcommand{\Kautz}{\mathrm K}
\newcommand{\Top}{\operatorname{Top}}
\newcommand{\cl}{\operatorname{cl}}
\newcommand{\Dspec}{\mathcal D}

\title{Symmetry premiums, cyclic-closure coverage, and\\
defect-spectrum persistence in Kautz digraphs}
\author{Oleksiy Babanskyy}
\date{}

\begin{document}
\pagestyle{plain}

\maketitle

\begin{abstract}
Let an involution of type \(2^m1^f\), \(m\geq1\), act on length-three
Kautz words by letterwise involution followed by reversal.  We prove that the
minimum invariant feedback vertex number is the ordinary value plus \(m^2\),
with surcharge \(m+m(m-1)\) split into local orbit forcing and a global
acyclicity cost.  Every invariant optimum is the full cyclic closure of an
ordinary optimum exactly when \(m=1\) in the nonidentity domain; the identity
case is stated separately.  For the order-cone defect spectrum we prove the
monotonicity
\[
  \Dspec_{r,0}\subseteq\Dspec_{m,f}
  \qquad(m\geq r\geq1,\ f\geq0).
\]
An exact finite census at \((m,f)=(3,0)\), accompanied here by one explicit
certificate for each value, therefore implies that defects \(0,1,2,3,4\)
are all attained for every \(m\geq3\) and \(f\geq0\).  This last conclusion
is attainability only: no equality of later spectra, general upper bound four,
or interval theorem is asserted.  More generally, defect is the transversal
number of the clutter of minimal infeasible bad-cell constraints.  Forced
cells and a matching of pairwise incompatibilities give a sound lower
certificate, while a four-letter rank-three example shows that this
certificate need not be exact.  At the six-letter base, a ten-relation
precedence core forces four bad triangle cells by a \(2+1+1\) argument and
occurs in exactly 24 of the 1,759 symmetry classes in the census.
\end{abstract}

\noindent\textbf{Keywords:} Kautz digraphs; feedback vertex sets;
reverse-complement symmetry; cyclic closure; defect spectra; exact census.

\medskip
\noindent\textbf{MSC 2020:} 05C20 (primary); 05C69 (secondary).

\medskip
\noindent\textbf{Reproducibility package:}
\url{https://github.com/aconsciousfractal/FCIG-Symmetry-Premiums-Cyclic-Closure-Coverage-and-Defect-Spectrum-Persistence-in-Kautz-Digraphs}

\section{Introduction}
\label{sec:introduction}

Feedback vertex sets in Kautz digraphs admit a useful line-digraph
description: a vertex deletion problem on length-three words becomes a
feedback arc problem on length-two words.  We study this problem under the
reverse-complement symmetry obtained by reversing a word and applying a
letter involution.  Three questions then separate naturally: what is the
exact cost of invariance, when are all invariant optima generated from
ordinary optima by cyclic closure, and which losses occur inside the
topological-order cone of an invariant optimum?  Here the cyclic closure of
a set \(F\) means \(\cl_{\Theta}(F)=F\cup\Theta_\sigma(F)\).

Our main contributions are as follows.
\begin{enumerate}[leftmargin=*]
\item For an involution of type \(2^m1^f\), the invariant optimum exceeds the
ordinary feedback vertex number by exactly \(m^2\).  The surcharge splits as
\(m+m(m-1)\) into a local orbit-forcing term and a global acyclicity term.
\item Complete cyclic-closure coverage holds exactly for one moving
transposition block; when \(m\geq2\), closure-generated and non-closure-
generated invariant optima coexist.
\item The defect spectrum is monotone under adjoining moving blocks and fixed
letters.  A certified six-letter base therefore yields persistent attained
defects \(0,1,2,3,4\) for every \(m\geq3\).
\item Defect equals the transversal number of an exact obstruction clutter.
Its singleton and pair members yield a forced-cell-plus-matching lower bound,
which a rank-three four-letter obstruction shows need not always be tight.
\item A ten-relation order core gives a short human explanation for a family
of extremal defect-four objects in the six-letter census.
\end{enumerate}

The ordinary value and cycle-cell substrate are known
\cite{XuWuHuangYang2007,XuYinZhangHuang2014}.  The feedback-arc/order and
line-digraph correspondences are classical
\cite{Younger1963,CheungKuh1974}.  We use only the elementary
antisymmetrization device behind Section~9 of Goldberg--Karzanov
\cite{GoldbergKarzanov2004}; no Kautz-specific theorem is imported from that
paper.  Adjacent work on symmetric decycling sets in binary de Bruijn graphs
uses reverse complementation to define the \(k\)-nonical space
\cite{MarcaisElderKingsford2024}.  Its graph family, alphabet, object and
optimization target differ from ours; we import no result from it and infer no
priority from that comparison.  Earlier work by the present author proves an
existence-and-parity criterion for literal reverse-complement-invariant
ordinary minimum decycling sets in \(q\)-ary de Bruijn graphs
\cite{Babanskyy2026Parity}.  The present paper instead treats length-three
Kautz digraphs, allows an invariance premium over the ordinary optimum, and
studies closure coverage and order-cone defect; no theorem is transferred
between the two graph families.  The full texts of two potentially relevant
Kautz sources
\cite{Yin2015Thesis,ZhangXuYangYin2016} could not be obtained.  We therefore
make no claim of bibliographic priority or source completeness.

\section{Ordinary Kautz substrate and notation}
\label{sec:ordinary}

Let \(A\) be an alphabet of size \(q\geq2\).  The vertices of
\(\Kautz(q-1,k)\) are the words \(x_1\cdots x_k\) with adjacent letters
distinct, and an arc shifts one position to the left.  Put
\[
 H=\Kautz(q-1,2),\qquad \Kautz(q-1,3)=L(H).
\]
Thus a word \(xyz\) is also the predecessor arc \(xy\to yz\) of \(H\).
Write
\[
 S_{q-1}=\frac{q(q-1)(2q-1)}6,
 \qquad
 \beta_q=\binom q2+4\binom q3.
\]

\begin{proposition}[Ordinary substrate; known background]
\label{prop:ordinary-substrate}
The minimum feedback vertex number of \(\Kautz(q-1,3)\) is \(S_{q-1}\),
and the maximum number of forward predecessor arcs in a total order of
\(V(H)\) is \(\beta_q\).
\end{proposition}

\begin{proof}
The arcs of \(H\) partition into one directed digon for each unordered pair
of letters and two oriented directed triangles for each unordered triple.
Every feedback arc set meets every cell, giving
\(\binom q2+2\binom q3=S_{q-1}\).  Given a total order on \(A\), delete the
strict middle peaks \(xyz\) with \(y>x,y>z\).  Their number is
\(\sum_{i=0}^{q-1}i^2=S_{q-1}\), and the standard descending-then-ascending
order of the pairs \(xy\) makes the complement acyclic.  Finally
\[
 q(q-1)^2-S_{q-1}=\binom q2+4\binom q3=\beta_q.
\]
The result and construction are cited background; the proof fixes our
normalization and the line-digraph crosswalk.
\end{proof}

\section{Symmetry, order cones, and defect}
\label{sec:setup}

Let \(\sigma\) be an involution of \(A\), of cycle type \(2^m1^f\), and put
\[
 \Theta_\sigma(xyz)=\sigma(z)\sigma(y)\sigma(x),
 \qquad j(xy)=\sigma(y)\sigma(x).
\]
On predecessor arcs the induced anti-action is
\(\rho(u,v)=(j(v),j(u))\).  Let \(\tau_\sigma\) be the minimum size of a
\(\Theta_\sigma\)-stable feedback vertex set and
\(\Delta_\sigma=\tau_\sigma-S_{q-1}\).

Let \(\mathcal C\) denote the digon-and-oriented-triangle partition above and
let \(G_\sigma=\langle\Theta_\sigma\rangle\) act on its cells.  For one
representative \(C\) of each cell orbit, put
\[
 H_C=\operatorname{Stab}_{G_\sigma}(C),
 \qquad s(C)=\min_{v\in C}|H_Cv|.
\]

\begin{proposition}[Invariant cell-transversal formula]
\label{prop:cell-orbit}
The minimum size of an invariant cell transversal is
\[
 L_\sigma=\sum_{[C]\in\mathcal C/G_\sigma}|G_\sigma C|s(C).
\]
Consequently
\[
 \ell_\sigma=L_\sigma-S_{q-1},\qquad
 g_\sigma=\tau_\sigma-L_\sigma,\qquad
 \Delta_\sigma=\ell_\sigma+g_\sigma.
\]
\end{proposition}

\begin{proof}
Let \(X\) be an invariant cell transversal and take \(v\in X\cap C\).
Invariance forces \(G_\sigma v\subseteq X\), while the cell partition gives
\(\operatorname{Stab}_{G_\sigma}(v)\leq H_C\).  Hence orbit--stabilizer gives
\[
 |G_\sigma v|
 =[G_\sigma:H_C]\,[H_C:\operatorname{Stab}_{G_\sigma}(v)]
 =|G_\sigma C|\,|H_Cv|.
\]
The inverse images of distinct cell orbits are disjoint, so summing gives the
lower bound.  Conversely, choose in every representative cell a vertex
attaining \(s(C)\) and take the union of its full \(G_\sigma\)-orbit.  This
invariant set meets every cell and attains the displayed sum.  Finally, the
cells are disjoint and their number is \(S_{q-1}\), whereas every feedback
set meets every cell.  Thus \(S_{q-1}\leq L_\sigma\leq\tau_\sigma\), proving
the decomposition and nonnegativity.
\end{proof}

Let \(P\) be a maximum \(\rho\)-stable acyclic predecessor arc set.  For
\(\omega\in\Top(P)\), write \(A_\omega\) for all ambient arcs of \(H\)
directed forward by \(\omega\), and define
\[
 \delta_H(P)=\beta_q-
 \max_{\omega\in\Top(P)}|A_\omega|.                 \tag{3.1}
\]
Let \(b(\omega)\) count the oriented-triangle cells having only one forward
arc in \(\omega\).

\begin{lemma}[Bad-triangle formula]
\label{lem:bad-triangle}
For every total order \(\omega\),
\(|A_\omega|=\beta_q-b(\omega)\); hence
\(\delta_H(P)=\min_{\omega\in\Top(P)}b(\omega)\).
\end{lemma}

\begin{proof}
Every digon has one forward arc.  Every directed triangle has one or two;
if \(b\) of the \(2\binom q3\) triangles have one, the score is
\(\binom q2+b+2(2\binom q3-b)=\beta_q-b\).
\end{proof}

For an ordinary feedback set \(F\), put
\(\cl_\Theta(F)=F\cup\Theta_\sigma(F)\).

\begin{lemma}[Order-cone dictionary; standard-derived]
\label{lem:order-cone}
The invariant minimum complementary to \(P\) is a full cyclic closure of an
ordinary minimum if and only if \(\delta_H(P)=0\).
\end{lemma}

\begin{proof}
Every acyclic arc set lies in the forward set of a topological order, and it
lies in an ordinary maximum precisely when such an order has score
\(\beta_q\).  If \(P\subseteq A_\omega\) with that score, stability gives
\(P\subseteq A_\omega\cap\rho(A_\omega)\); the intersection is stable and
acyclic, so maximality forces equality.  Complementation gives the full
closure of the ordinary minimum \(E(H)\setminus A_\omega\).  Conversely, if
the complement of \(P\) is such a closure, then \(P\) is contained in the
complementary ordinary maximum and therefore has a topological order of
score \(\beta_q\).
\end{proof}

\section{Exact symmetry premium and its split}
\label{sec:premium}

For the fixed-point-free case write the transposition blocks as
\(B_i=\{a_i,b_i\}\), \(0\leq i<m\), and \(\bar a_i=b_i\).

\begin{theorem}[Exact premium and local/global split]
\label{thm:premium}
If \(\sigma\) has type \(2^m\), \(m\geq1\), on \(q=2m\) letters, then
\[
 \tau_\sigma(\Kautz(2m-1,3))=S_{2m-1}+m^2,
 \quad \Delta_\sigma=m^2,
 \quad \ell_\sigma=m,
 \quad g_\sigma=m(m-1).
\]
\end{theorem}

\begin{proof}
By Lemma~\ref{lem:potential-completion} below, a maximum stable predecessor
DAG is \(I(h)=\{u\to v:h(u)<h(v)\}\) for an antisymmetric potential
\(h\circ j=-h\).  The fixed predecessor vertices \(a_ib_i,b_ia_i\) have
level zero.  Classify an arc \(xy\to yz\) by the set of transposition blocks
meeting \(\{x,y,z\}\).  Its block support has size one, two, or three, and
the following table is an exhaustive disjoint partition.  The counts and
bounds are for each fixed choice of block support.
\begin{center}
\begin{tabularx}{\textwidth}{@{}cXcc@{}}
\toprule
support & pieces & number & bound per piece\\
\midrule
one block & within-block digon & 1 & 0\\
two blocks & fixed--free two-arc pairs & 8 & 1\\
           & cross-block digons & 4 & 1\\
           & non-digon directed four-cycles & 2 & 3\\
three blocks & directed triangles & 16 & 2\\
\bottomrule
\end{tabularx}
\end{center}
Indeed, in one block both predecessor vertices are fixed and tied at zero.
For two blocks there are eight nonfixed predecessor vertices; each supports
one incoming arc from a fixed vertex and one outgoing arc to a fixed vertex,
so at most one of the two can increase.  The remaining arcs are exactly four
digons and two directed four-cycles.  For three blocks they are exactly the
two oriented triangles on each of the eight choices of one letter per block.
An increasing potential selects at most one arc of a digon, three of a
directed four-cycle, and two of a directed triangle.  Therefore
\[
 |I(h)|\leq B_m:=18\binom m2+32\binom m3
 =4m(m-1)+10\binom m2+32\binom m3
 =\frac{m(m-1)(16m-5)}3.                           \tag{4.1}
\]

This bound is sharp.  Put \(\lambda_i=4^{m-1-i}\), set
\(c_{ij}=2\lambda_j\) for \(i<j\), \(c_{ji}=-c_{ij}\), and identify
\(i^-=a_i,i^+=b_i\).  For \(i\ne j\) define
\[
 h(i^\varepsilon j^\delta)
 =\varepsilon\lambda_i+\delta\lambda_j+c_{ij}
 \qquad(\varepsilon,\delta\in\{-1,+1\}).             \tag{4.2}
\]
Complete the definition by
\(h(i^-i^+)=h(i^+i^-)=0\).  Since
\(j(i^\varepsilon j^\delta)=j^{-\delta}i^{-\varepsilon}\) and
\(c_{ji}=-c_{ij}\), this is antisymmetric.  If \(i<j\), the sign of
\(h(i^\varepsilon j^\delta)\) is \(\varepsilon\), since
\(\lambda_i\geq4\lambda_j\).  Thus all eight fixed--free pairs in that
two-block sector attain their bound.

Fix \(i<j\) and let
\(\varepsilon,\delta,\eta,\zeta\in\{-1,+1\}\).  On the two alternating
arc families
\[
 i^\varepsilon j^\delta\to j^\delta i^\eta,
 \qquad
 j^\delta i^\varepsilon\to i^\varepsilon j^\zeta,
\]
the potential differences are respectively
\[
 (\eta-\varepsilon)\lambda_i-4\lambda_j,
 \qquad (\zeta-\delta)\lambda_j+4\lambda_j;
\]
the first is positive exactly for
\((\varepsilon,\eta)=(-1,+1)\), giving two arcs as \(\delta\) varies,
while the second belongs to \(\{2\lambda_j,4\lambda_j,6\lambda_j\}\) and
is positive for all eight choices.  Hence the four digons and two
four-cycles contribute exactly \(2+8=10\).

Finally fix \(i<j<k\) and
\(\varepsilon,\delta,\gamma\in\{-1,+1\}\).  Around the directed triangle
\[
 i^\varepsilon j^\delta\to j^\delta k^\gamma
 \to k^\gamma i^\varepsilon\to i^\varepsilon j^\delta
\]
the three potential differences are
\[
 -\varepsilon\lambda_i-2\lambda_j+(\gamma+2)\lambda_k,
 \quad \varepsilon\lambda_i-\delta\lambda_j-4\lambda_k,
 \quad (\delta+2)\lambda_j+(2-\gamma)\lambda_k.
\]
Their signs are \((-\varepsilon,\varepsilon,+)\), since
\(\lambda_i\geq4\lambda_j\geq16\lambda_k\); exactly two arcs increase in
each of these eight triangles.  Antisymmetry makes \(I(h)\) \(\rho\)-stable;
the other eight oriented triangles are their \(\rho\)-images as the signs
vary, and have the same count.  Thus every row
of the exhaustive sector table is attained.

Since \(H\) has \(2m(2m-1)^2\) arcs, complementation gives
\[
 \tau_\sigma=2m(2m-1)^2-B_m=S_{2m-1}+m^2.
\]
For the local term, precisely the \(m\) within-block digon cells are fixed
with their two vertices exchanged, forcing two rather than one selected
vertex.  Every other cell orbit has trivial internal stabilizer and attains
the ordinary cell cost.  Proposition~\ref{prop:cell-orbit} gives
\(L_\sigma=S_{2m-1}+m\), and hence
\(\ell_\sigma=m\) and \(g_\sigma=m^2-m\).
\end{proof}

\section{Fixed letters and zero/one constructions}
\label{sec:constructions}

\begin{lemma}[Exact fixed-letter transport]
\label{lem:fixed-letter}
Adjoining one \(\sigma\)-fixed letter \(t\) to an alphabet \(B\) of size
\(n\) increases the ordinary and invariant feedback optima by \(n^2\),
preserves the symmetry premium and its local/global split, and sends a
canonical lifted maximum predecessor DAG to one of exactly the same defect.
\end{lemma}

\begin{proof}
For a feedback set \(M_B\), put
\[
 M_{B+t}=M_B\cup X_t,\qquad X_t=\{xtz:x,z\in B\}.       \tag{5.1}
\]
Extend a topological order of the old residual by placing all \(ta\) before
all old pairs and all \(at\) after them.  The retained new words are forward,
so (5.1) is a feedback set.  The \(n\) new digons and
\(2\binom n2\) new oriented triangles are disjoint and total \(n^2\); every
feedback set must meet all of them.  This proves the ordinary increment.
The set \(X_t\) is invariant, and restriction of an invariant feedback set is
invariant, giving the same invariant increment.  Its unique middle-\(t\)
choice meets each new cell orbit minimally, so the split is preserved.

Let \(P^\uparrow\) be the residual of (5.1).  Restricting a topological order
of \(P^\uparrow\) to \(B\) shows
\(\delta_{B+t}(P^\uparrow)\geq\delta_B(P)\), because new cells contribute
at most \(\beta_{n+1}-\beta_n\).  The displayed extended order attains the
maximum in every new cell, giving the reverse inequality.  Thus the defect
is preserved, and the construction iterates.
\end{proof}

\begin{corollary}[Premium with fixed letters]
\label{cor:premium-general}
For type \(2^m1^f\), \(m\geq1,f\geq0,q=2m+f\),
\[
 \tau_\sigma(\Kautz(q-1,3))=S_{q-1}+m^2,
 \quad(\ell_\sigma,g_\sigma)=(m,m(m-1)).
\]
\end{corollary}

\begin{proof}
Iterate Lemma~\ref{lem:fixed-letter} from Theorem~\ref{thm:premium}.
\end{proof}

\begin{proposition}[Four-letter seeds]
\label{prop:four-letter}
For \(A=\{0,1,2,3\}\) and \(\sigma=(01)(23)\), maximum invariant
predecessor DAGs of defects zero and one both exist.
\end{proposition}

\begin{proof}
Here \(\beta_4=22\) and every invariant minimum has size \(18\).  For defect
zero, take the ordinary minimum
\[
 M=\{010,012,132,202,210,212,232,301,302,303,310,312,313,320\}
\]
and its closure
\[
 \begin{aligned}
 N_0=M\cup\Theta(M)=\{&010,012,101,102,103,132,202,210,212,\\
                       &232,301,302,303,310,312,313,320,323\}.
 \end{aligned}
\]
The complement of \(M\) is forward in
\(10<02<12<20<01<03<32<21<13<23<30<31\), while the complement of
\(N_0\) is forward in
\(02<10<12<20<01<03<23<32<21<13<30<31\).  Hence both are optimal and
the residual of \(N_0\) has defect zero by Lemma~\ref{lem:order-cone}.

For defect one, let
\[
 N_1=\{010,013,021,030,031,032,101,102,121,123,201,202,
 230,232,310,313,321,323\}.
\]
It is invariant, and its complement is forward in
\[
 21<10<13<30<01<02<23<31<12<32<20<03.
\]
This order scores \(21\), with unique bad cell
\(\{031,310,103\}\).  Defect zero would force an ordinary minimum
\(M\subseteq N_1\) whose cyclic closure is \(N_1\).  Such an \(M\) contains
exactly one word from each of the fourteen cycle cells.  Ten singleton
intersections force
\[
 202,030,121,313,201,013,230,032,123,321.
\]
The remaining digons force one of \(010,101\) and one of \(232,323\).
The final two triangle choices must be either \(\{021,310\}\) or
\(\{102,031\}\).  In the first case the residual contains
\[
 312\to120\to203\to031\to312,
\]
and in the second it contains
\[
 213\to130\to302\to021\to213.
\]
Neither candidate is a feedback vertex set, independently of the two digon
choices.  Thus the defect is positive and therefore exactly one.
\end{proof}

\begin{theorem}[Uniform zero/one construction]
\label{thm:zero-one}
For every \(m\geq1,f\geq0\), defect zero is attained.  For every
\(m\geq2,f\geq0\), defect one is attained.
\end{theorem}

\begin{proof}
For \(m=1\), the empty moved-core is the unique maximum and has defect zero.
Lemma~\ref{lem:potential-completion} and
Proposition~\ref{prop:sharp-core} represent this seed and both four-letter
seeds of Proposition~\ref{prop:four-letter} as \(I(h)\) for sharp
antisymmetric potentials.
Lemmas~\ref{lem:dominant-extension} and \ref{lem:defect-extension} below
extend a sharp core by one moving block while preserving its exact defect.
They therefore propagate the zero seed from \(m=1\), and the defect-one seed
of Proposition~\ref{prop:four-letter} from \(m=2\), to every larger \(m\).
Lemma~\ref{lem:fixed-letter} then propagates both families over all \(f\).
The cited lemmas are independently proved in Section~\ref{sec:extension}, so
this forward reference is not circular.
\end{proof}

\section{Fixed-digon ceiling and closure coverage}
\label{sec:coverage}

\begin{proposition}[Fixed-digon defect ceiling]
\label{prop:ceiling}
Every maximum \(\rho\)-stable predecessor DAG \(P\) of type \(2^m1^f\),
\(m\geq1\), satisfies
\[
 0\leq\delta_H(P)\leq m(m-1).
\]
\end{proposition}

\begin{proof}
For each moving block, the vertices \(a_ib_i,b_ia_i\) are fixed by \(j\),
while the two arcs between them are exchanged by \(\rho\).  A stable DAG
contains neither.  Every topological order recovers exactly one forward arc
from each such omitted digon.  By Corollary~\ref{cor:premium-general},
\(|P|=\beta_q-m^2\), so every \(\omega\in\Top(P)\) has
\(|A_\omega|\geq\beta_q-m^2+m\).  Subtracting the maximum score from
\(\beta_q\) proves the bound.
\end{proof}

\begin{theorem}[Cyclic-closure coverage]
\label{thm:coverage}
For every nonidentity involution of type \(2^m1^f\), \(m\geq1,f\geq0\):
every maximum invariant predecessor DAG has defect in
\([0,m(m-1)]\); a zero-defect maximum exists; an exact defect-one maximum
exists for \(m\geq2\); and every invariant minimum is a full cyclic closure
of an ordinary minimum if and only if \(m=1\).
\end{theorem}

\begin{proof}
The bound is Proposition~\ref{prop:ceiling}, and the two existence assertions
are Theorem~\ref{thm:zero-one}.  If \(m=1\), the bound forces every maximum
to have defect zero, and Lemma~\ref{lem:order-cone} gives complete coverage.
If \(m\geq2\), the exact defect-one construction supplies an invariant
minimum that is not closure-generated.  The simultaneous zero-defect
construction shows that existence of a closure-generated optimum and failure
of complete coverage coexist.
\end{proof}

\begin{proposition}[Identity boundary]
\label{prop:identity}
If \(\sigma=\mathrm{id}\), \(q=f\geq2\), then
\(\tau_\sigma=S_{q-1}\), the surcharge is zero, and every invariant minimum
is the full cyclic closure of an ordinary minimum.
\end{proposition}

\begin{proof}
The strict middle-peak ordinary minimum from
Proposition~\ref{prop:ordinary-substrate} is invariant under reversal, so the
ordinary and invariant values agree.  Any invariant minimum is itself an
ordinary minimum and equals its own cyclic closure.
\end{proof}

\section{Sharp potentials and support extension}
\label{sec:extension}

For a potential \(h:V(H)\to\mathbb R\), write
\(I(h)=\{u\to v\in E(H):h(u)<h(v)\}\).

\begin{lemma}[Antisymmetric potential completion]
\label{lem:potential-completion}
Let a finite digraph have a vertex involution \(j\), with arc anti-action
\(\rho(u,v)=(j(v),j(u))\).  If \(P\) is a \(\rho\)-stable acyclic arc set,
then an antisymmetric potential \(h\circ j=-h\) satisfies \(P\subseteq I(h)\).
If \(P\) has maximum cardinality among stable acyclic arc sets, then
\(P=I(h)\).
\end{lemma}

\begin{proof}
Choose a strictly increasing topological rank \(t\) of \(P\) and put
\(h(v)=t(v)-t(j(v))\).  Then \(h(j(v))=-h(v)\).  If \(u\to v\in P\),
stability also gives \(j(v)\to j(u)\in P\), and
\[
 h(v)-h(u)=[t(v)-t(u)]+[t(j(u))-t(j(v))]>0.
\]
Thus \(P\subseteq I(h)\).  The set \(I(h)\) is acyclic because \(h\) strictly
increases on it, and it is stable by antisymmetry.  Maximal cardinality now
forces equality.
\end{proof}

\begin{proposition}[Sharp-core localization]
\label{prop:sharp-core}
In fixed-point-free type \(2^m\), every maximum stable predecessor DAG is
\(I(h)\) for a sharp antisymmetric potential attaining \(B_m\).  Its
restriction to any subset of moving blocks is maximum.  Moreover, if
\(P_{ij}\) is its two-block restriction, then
\[
 \delta_H(P)\geq\sum_{i<j}\delta(P_{ij}).              \tag{7.1}
\]
\end{proposition}

\begin{proof}
Lemma~\ref{lem:potential-completion} gives \(P=I(h)\).  Theorem
\ref{thm:premium} and (4.1) give \(|P|=B_m\).  Since (4.1) is the sum of
disjoint sector bounds, equality forces saturation of every sector; summing
the sectors supported on any chosen block subset gives its exact smaller
bound.  Thus every restriction is maximum.  Finally, by
Lemma~\ref{lem:bad-triangle}, defect counts bad oriented triangles.  The
bad triangle cells supported on different unordered block pairs are disjoint
(shared within-block digons have zero defect), and every global topological
order restricts to one of each \(P_{ij}\).  Summing their local lower bounds
proves (7.1).  Equality in (7.1) is not asserted.
\end{proof}

\begin{lemma}[Sharp dominant moving-block extension]
\label{lem:dominant-extension}
Let \(h\) be an antisymmetric potential on \(r\geq1\) moving blocks with
\(|I(h)|=B_r\).  Put \(H_0=\max_v|h(v)|\), \(L=H_0+4\), and adjoin the
block \(\{a,b\}\).  If \(\varepsilon_x\in\{-1,+1\}\) is the sign of an old
letter in its block, preserve all old levels and set
\[
\begin{array}{rclrcl}
 \widetilde h(ax)&=-L+\varepsilon_x+2,&
 \widetilde h(xa)&=-L+\varepsilon_x-2,\\
 \widetilde h(bx)&= L+\varepsilon_x+2,&
 \widetilde h(xb)&= L+\varepsilon_x-2,\\
 \widetilde h(ab)&=0,&\widetilde h(ba)&=0.
\end{array}                                             \tag{7.2}
\]
Then \(\widetilde h\) is antisymmetric and
\(|I(\widetilde h)|=B_{r+1}\).
\end{lemma}

\begin{proof}
Write the new letters as \(n^-=a,n^+=b\), and write an old letter as
\(x^\delta\), where \(\delta\in\{-1,+1\}\) records its sign in its moving
block.  Then (7.2) is equivalently
\[
 \widetilde h(n^\varepsilon x^\delta)=\varepsilon L+\delta+2,
 \qquad
 \widetilde h(x^\delta n^\varepsilon)=\varepsilon L+\delta-2.
\]
These formulas and
\(j(n^\varepsilon x^\delta)=x^{-\delta}n^{-\varepsilon}\) prove
antisymmetry directly.  Every new nonfixed level has absolute value at least
\(H_0+1\), so each of the eight fixed--free two-arc pairs belonging to a
new/old block pair contributes exactly one arc.

For \(\varepsilon,\zeta,\delta,\gamma\in\{-1,+1\}\), the two alternating
families have differences
\[
 (\zeta-\varepsilon)L-4,
 \qquad \gamma-\delta+4\in\{2,4,6\};
\]
the first is positive exactly when
\((\varepsilon,\zeta)=(-1,+1)\), giving two arcs as \(\delta\) varies, and
the second is positive for all eight choices.  Hence the alternating part
attains ten, just as in the two-block row of the sector table.

Now choose two distinct old letters \(x^\delta,y^\gamma\) from different
old blocks.  On the directed triangle
\(n^\varepsilon x^\delta\to x^\delta y^\gamma
\to y^\gamma n^\varepsilon\to n^\varepsilon x^\delta\), its cyclic levels
have the form
\[
 p=\varepsilon L+\delta+2,
 \quad q\in[-H_0,H_0],
 \quad r=\varepsilon L+\gamma-2.
\]
If \(\varepsilon=-1\), then \(r<p<q\); if \(\varepsilon=+1\), then
\(q<r<p\).  Exactly two cyclic arcs increase in either case, and the reverse
orientation is covered by interchanging \(x^\delta\) and \(y^\gamma\).
Thus all sixteen triangle cells for the new block and each old block pair
attain two.  The within-new-block digon remains tied at zero, old sectors
remain sharp, and the total increment is
\[
 18r+32\binom r2=B_{r+1}-B_r.
\]
Consequently \(|I(\widetilde h)|=B_{r+1}\), as required.
\end{proof}

\begin{lemma}[Exact moving-block defect preservation]
\label{lem:defect-extension}
For \(P=I(h)\) and \(\widetilde P=I(\widetilde h)\) as in
Lemma~\ref{lem:dominant-extension},
\[
 \delta_{A\cup\{a,b\}}(\widetilde P)=\delta_A(P).       \tag{7.3}
\]
\end{lemma}

\begin{proof}
Let \(\omega\in\Top(P)\).  Extend it by the seven bands
\[
 ba;\quad xa\ (x\in A);\quad ab;\quad ax\ (x\in A);\quad
 \omega;\quad xb\ (x\in A);\quad bx\ (x\in A),         \tag{7.4}
\]
using any fixed old-letter order inside homogeneous bands.  Directly from
(7.2), every selected new arc goes forward in (7.4): arcs from \(xa\) to
\(ab,ay\), from \(ab\) to \(bx\), from \(ax\) to old vertices or \(xb\),
from old vertices to \(yb\), and from \(xb\) to \(by\) all respect the band
order; the possible reverse-band arcs have decreasing or tied potential.
The new triangle restrictions are
\[
 \begin{aligned}
 xa&<ab<bx, & ba&<ax<xb, & ya&<ax<xy,\\
 xa&<ay<yx, & xy&<yb<bx, & yx&<xb<by.
 \end{aligned}
\]
Thus every new digon scores one and every new triangle scores two, while old
cells retain their score under \(\omega\).  Taking a maximizing \(\omega\)
gives \(\delta(\widetilde P)\leq\delta(P)\).

Conversely, for any alphabet inclusion \(A\subseteq A'\), restriction of a
topological order of \(P'\) is topological for \(P=P'|_A\), and new cells
score at most \(\beta_{|A'|}-\beta_{|A|}\).  Hence
\[
 \max s_{A'}\leq\max s_A+\beta_{|A'|}-\beta_{|A|},
 \qquad \delta_{A'}(P')\geq\delta_A(P).
\]
Apply this to \(\widetilde P|_A=P\) to obtain the reverse inequality in
(7.3).
\end{proof}

\section{Defect-spectrum monotonicity}
\label{sec:monotonicity}

Let \(\Dspec_{m,f}\) be the set of defects attained by maximum stable
predecessor DAGs for type \(2^m1^f\).

\begin{theorem}[Support-extension monotonicity]
\label{thm:monotonicity}
For every \(m\geq r\geq1\) and \(f\geq0\),
\[
 \Dspec_{r,0}\subseteq\Dspec_{m,f}.
\]
\end{theorem}

\begin{proof}
Fix \(d\in\Dspec_{r,0}\) and a maximum stable DAG \(P\) of defect \(d\).
Lemma~\ref{lem:potential-completion} writes \(P=I(h)\); Proposition
\ref{prop:sharp-core} says that \(h\) is sharp.  Lemma
\ref{lem:dominant-extension} may therefore add one moving block, and Lemma
\ref{lem:defect-extension} preserves the defect exactly.  Iteration gives a
maximum of defect \(d\) at \((m,0)\).  Finally, iterate the canonical
fixed-letter lift of Lemma~\ref{lem:fixed-letter} \(f\) times.  It also
preserves maximality and defect exactly, proving the inclusion.
\end{proof}

\begin{corollary}[Five persistent attained defects]
\label{cor:five-defects}
For every \(m\geq3,f\geq0\),
\[
 \{0,1,2,3,4\}\subseteq\Dspec_{m,f}.
\]
\end{corollary}

\begin{proof}
The exact six-letter census established below gives
\(\Dspec_{3,0}=\{0,1,2,3,4\}\).  Apply
Theorem~\ref{thm:monotonicity} with \(r=3\).  Only the five existence
statements are used here; the exhaustive upper-bound part of the census is not
transported.
\end{proof}

\begin{remark}[Logical boundary]
The theorem does not say that \(\Dspec_{m,f}\) equals the five displayed
values, that its maximum is four, or that it is an interval.  It transports
neither realization-cone facets nor adjacency, connectivity, extrema
reachability, or Lipschitz properties.
\end{remark}

\section{Five explicit finite seeds and the computational boundary}
\label{sec:seeds}

At \(q=6\), \(\sigma=(01)(23)(45)\), the 150 Kautz words form 75
two-element \(\rho\)-orbits in lexicographic order.  In this finite census, a
\emph{labelled maximum mask} is the 75-bit incidence vector of a maximum
\(\rho\)-stable acyclic predecessor-arc set; it has weight 43 and therefore
selects 86 arcs.  A \emph{source triple} is the ordered triple of local
four-letter pair-core states on block pairs \((01),(02),(12)\).  Each local
state lies in a canonically ordered 36-element list, so indices
\((i_{01},i_{02},i_{12})\) have radix-36 code
\(36^2i_{01}+36i_{02}+i_{12}\).  A \emph{closure--centralizer work unit} is
an equivalence class of labelled maximum masks whose transitive closures
agree after relabelling by the full centralizer of \(\sigma\).  A
\emph{mask key} is the zero-padded 19-digit hexadecimal encoding of the
75-bit vector.  In the last table column, \emph{relations/covers} means the
cardinalities of the transitive closure and transitive reduction,
respectively.  The following deterministic certificates give one witness
for each defect.

\begin{center}
\small
\begin{tabular}{cclrr}
\toprule
defect & score & mask key & source triple & relations/covers\\
\midrule
0 & 95 & \texttt{06d87f4fef6017ea3d6} & 5961  & 334/40\\
1 & 94 & \texttt{25d86fc9dfe7f1027e4} & 25902 & 332/42\\
2 & 93 & \texttt{06dc074fcf67d3ea3f4} & 26734 & 350/38\\
3 & 92 & \texttt{24dc67cbdfecf3823e2} & 12907 & 334/56\\
4 & 91 & \texttt{25d82fcdcffff2007f0} & 38610 & 362/48\\
\bottomrule
\end{tabular}
\end{center}

The associated optimal predecessor orders and their bad cells are:
\begin{enumerate}[label={\(d=\arabic*:\)},start=0,leftmargin=*]
\item {\ttfamily\footnotesize 01 03 04 31 14 41 10 12 13 15 53 34 42 23 43 30 32 25\\
35 45 52 20 24 54 40 05 50 02 21 51}\par No bad cell.
\item {\ttfamily\footnotesize 01 04 31 42 52 21 14 24 41 10 12 13 15 23 25 45 54 40\\
43 30 32 20 05 50 02 03 53 34 35 51}\par Bad cell \(T^-_{235}\).
\item {\ttfamily\footnotesize 01 31 12 52 24 41 42 21 15 23 25 51 10 13 14 32 20 45\\
50 04 54 40 03 05 43 35 53 30 02 34}\par Bad cells \(T^-_{234},T^-_{235}\).
\item {\ttfamily\footnotesize 12 35 40 02 52 20 01 05 21 23 25 51 10 14 45 50 03 04\\
53 32 24 54 41 13 43 30 31 15 34 42}\par Bad cells
\(T^+_{015},T^-_{123},T^+_{245}\).
\item {\ttfamily\footnotesize 12 50 02 20 24 40 01 03 04 32 34 41 10 05 42 23 25 45\\
54 43 35 51 52 21 13 15 53 30 31 14}\par Bad cells
\(T^-_{012},T^-_{013},T^-_{145},T^-_{234}\).
\end{enumerate}
Each order respects the complete canonical precedence relation encoded by
its mask, and Lemma~\ref{lem:bad-triangle} gives the displayed defect.
The complete transitive closures, reductions, centralizer transports, opposed
dynamic-programming certificates, and hashes are stored in a machine-readable
certificate bundle.

The existence list alone does not prove that no larger defect occurs.  The
population is reconstructed independently of the stored labels and mask table
from the combinatorial definition, using the frozen quotient-native Route B
kernel.  This is a label-independent replay, not a second implementation of
that kernel.  The portable enumerator first derives
the 36 local four-letter states, forms all \(36^3=46{,}656\) source triples,
quotients them into 987 centralizer orbits, and visits exactly 54,078 nodes in
the complete leaf-canonical triangle search.  It obtains 1,759 canonical
representative masks and expands them under the 48-element centralizer to
exactly 83,736 distinct labelled maximum masks.  Only after this enumeration
does it open the tracked table and verify exact set equality and digest.

The separate portable exhaustive optimization then covers all 83,736
labelled maximum masks and all 1,759 closure--centralizer work units.  For
every work unit it checks
opposed prefix and suffix dynamic-programming tables state by state against
their Bellman recurrences, verifies that both routes give the same optimum,
and checks a concrete optimal topological order.  Its
labelled-mask multiplicities for defects \(0,\ldots,4\) are
\[
 (20976,26640,19896,14208,2016),
\]
and its work-unit multiplicities are
\[
 (445,555,420,297,42).
\]
That computation proves the upper boundary only at \((m,f)=(3,0)\).  The five
certificates make the existence witnesses portable; they are not a
replacement for the exhaustive census.

\section{Exact obstruction clutters and lower certificates}
\label{sec:obstruction-clutter}

For distinct letters \(a<b<c\), let \(T^+_{abc}\) and \(T^-_{abc}\) denote
the two oriented triangle cells, with cyclic comparisons
\[
 \begin{array}{lll}
 T^+_{abc}:&ab<bc,&bc<ca,\quad ca<ab,\\
 T^-_{abc}:&ac<cb,&cb<ba,\quad ba<ac.
 \end{array}
\]
A total order makes either one or two comparisons in each cell true; the cell
is bad in the former case.  Let \(\mathcal T_q\) be the set of all these
oriented triangle cells.  In this section formula (3.1) defines
\(\delta_H(P)\) for any acyclic predecessor arc set, not necessarily a
maximum one.  For such a set \(P\) and
\(S\subseteq\mathcal T_q\), write \(\operatorname{Good}_P(S)\) when some
order in \(\Top(P)\) makes every cell of \(S\) good.  Define the obstruction
clutter
\[
 \mathcal H(P)=\min_{\subseteq}
 \{S\subseteq\mathcal T_q:\operatorname{Good}_P(S)\text{ fails}\},
                                                               \tag{10.1}
\]
where only inclusion-minimal members are retained.  These are cyclic-order
constraints in the sense of \cite{GalilMegiddo1977,Opatrny1979}; their
incidence vectors lie naturally in the linear-ordering-polytope setting of
\cite{GroetschelJuengerReinelt1985}.  The following identity is elementary
and does not import a result from those papers.

\begin{theorem}[Exact obstruction-clutter formula]
\label{thm:obstruction-clutter}
For every acyclic predecessor arc set \(P\),
\[
 \delta_H(P)=\tau(\mathcal H(P)),                         \tag{10.2}
\]
where the right-hand side is the transversal number of the clutter.
\end{theorem}

\begin{proof}
For \(\omega\in\Top(P)\), let \(B(\omega)\) be its set of bad oriented
triangle cells.  This set meets every member of \(\mathcal H(P)\), since a
disjoint obstruction would be good under \(\omega\).  Lemma~\ref{lem:bad-triangle}
therefore gives \(\delta_H(P)\geq\tau(\mathcal H(P))\).

Conversely, let \(B\) be a minimum transversal.  The set
\(\mathcal T_q\setminus B\) contains no member of \(\mathcal H(P)\).  If it
were infeasible, finiteness would give an inclusion-minimal infeasible subset,
contradicting that assertion.  Hence some \(\omega\in\Top(P)\) makes every
cell outside \(B\) good, so \(B(\omega)\subseteq B\).  The bad-triangle
formula yields the reverse inequality.
\end{proof}

Let
\[
 F(P)=\{T\in\mathcal T_q:\operatorname{Good}_P(\{T\})\text{ fails}\}.
\]
On \(\mathcal T_q\setminus F(P)\), form the incompatibility graph \(G(P)\)
by joining \(T\) and \(U\) when
\(\operatorname{Good}_P(\{T,U\})\) fails.

\begin{corollary}[Forced cells plus matching]
\label{cor:forced-matching}
For every acyclic \(P\),
\[
 \delta_H(P)\geq |F(P)|+\nu(G(P)).                       \tag{10.3}
\]
\end{corollary}

\begin{proof}
The singleton members of \(\mathcal H(P)\) are precisely the forced cells.
After deleting them, every edge of \(G(P)\) is a two-element member of the
clutter.  A matching gives pairwise disjoint such obstructions, so every
transversal contains all forced cells and at least one endpoint of each
matching edge.  Apply Theorem~\ref{thm:obstruction-clutter}.
\end{proof}

The lower certificate (10.3) is not always exact.  The following small
control prevents higher-rank obstructions from being silently discarded.

\begin{proposition}[A rank-three obstruction]
\label{prop:rank-three-obstruction}
On four letters, let \(P\) be the partial order generated by
\[
\begin{array}{lllll}
01<13,&03<20,&03<32,&10<31,&12<01,\\
20<23,&21<32,&31<20,&32<02,&32<30.
\end{array}                                             \tag{10.4}
\]
Then \(F(P)=\varnothing\), \(G(P)\) has no edge, but
\(\delta_H(P)=1\).  Its unique minimal obstruction is
\[
 \{T^+_{012},T^-_{023},T^-_{123}\}.
\]
Thus equality in (10.3) is false in general.
\end{proposition}

\begin{proof}
Direct enumeration gives 99,000 linear extensions of \(P\), 162 with exactly
one bad cell, and 176 distinct bad-cell sets.  It verifies that every singleton
and every pair of cells is simultaneously good, while every proper subset of
the displayed triple is feasible.  The triple itself has a short certificate:
if all three cells were good, the base relations leave only compatible cyclic
chains containing respectively
\[
 20<01,\qquad32<20,\qquad13<32.
\]
Together with \(01<13\), these give
\(20<01<13<32<20\), a contradiction.  Hence the displayed triple is the
unique minimal obstruction and Theorem~\ref{thm:obstruction-clutter} gives
\(\delta_H(P)=1\).
\end{proof}

\section{A six-letter four-defect obstruction}
\label{sec:obstruction}

We now apply the forced-cell-plus-matching certificate to a six-letter core.

\begin{lemma}[Ten-relation four-defect obstruction]
\label{lem:ten-relation-obstruction}
Every total order containing the precedence relations
\[
\begin{array}{lll}
02<10<21, & 03<10<31,\\
40<54<15, & 41<05,\\
24<32<53, & 25<43
\end{array}                                             \tag{11.1}
\]
has at least four bad cells among
\[
T^-_{012},T^-_{013},T^-_{045},T^-_{145},T^-_{234},T^-_{235}.
\]
\end{lemma}

\begin{proof}
The chain \(02<10<21\) gives the truth pattern \((1,0,0)\) in
\(T^-_{012}\), so this cell is bad.  Likewise \(03<10<31\) makes
\(T^-_{013}\) bad.

Suppose that both \(T^-_{045}\) and \(T^-_{145}\) were good.  Since
\(40<54\), goodness of the first cell would force \(40<05<54\).  Since
\(54<15\), goodness of the second would force \(54<41<15\).  Hence
\(05<54<41\), contrary to \(41<05\).  At least one of these two cells is
therefore bad.

Similarly, goodness of \(T^-_{234}\) would force \(24<43<32\), while
goodness of \(T^-_{235}\) would force \(32<25<53\).  Together they imply
\(43<32<25\), contrary to \(25<43\).  At least one of the last two cells is
bad.  The three disjoint contributions are consequently \(2+1+1=4\).
Equivalently, if \(P\) is the partial order generated by (11.1), the first
two cells lie in \(F(P)\), while
\(\{T^-_{045},T^-_{145}\}\) and
\(\{T^-_{234},T^-_{235}\}\) are disjoint edges of \(G(P)\);
Corollary~\ref{cor:forced-matching} gives the same lower bound.
\end{proof}

The argument is invariant under relabelling and remains valid after adding
vertices or further precedence relations, provided the displayed core is
retained.  It does not assert that the core occurs for every involution or in
every extremal order cone.

\begin{proposition}[Exact occurrence in the six-letter census]
\label{prop:obstruction-occurrence}
Let \(\sigma=(01)(23)(45)\).  Among the 1,759 closure-centralizer classes in
the six-letter census, exactly 24 contain an image of the core (11.1) under
\(C_{S_6}(\sigma)\).  Each of these 24 classes has defect four, each contains
a unique centralizer image of the core, and none of the remaining 1,735
classes contains such an image.  The 24 classes account for 1,152 labelled
masks.
\end{proposition}

\begin{proof}[Computer verification]
A separate checker reconstructs the 48-element centralizer, decodes every
closure, tests all centralizer images of (11.1), and independently checks all
1,152 labelled members against their canonical representatives.  It records
24 unique hits, no hit outside the selected classes, and no failed transport.
The commands and certificate boundary are given in
Appendix~\ref{app:reproducibility}.
\end{proof}

\begin{remark}[Minimality boundary]
Within this fixed six-cell decomposition, deleting any one of the ten
relations lowers the forced minimum in its component, so (11.1) is
inclusion-minimal for the displayed \(2+1+1\) proof.  It is not claimed to be
a globally smallest classifier of the 24 finite classes.
\end{remark}

\section{Finite diagnostics and open mechanism questions}
\label{sec:diagnostics}

We first fix the geometric terminology used in this section.  For a fixed
involution \(\sigma\), let
\[
 V^-_\sigma=\{h:A_q\to\mathbb R:h(jv)=-h(v)\},\qquad
 \ell_{xyz}(h)=h(yz)-h(xy).
\]
The forms are constant on the \(\rho\)-orbits of words.  After removing the
structural zero forms and identifying proportional forms, their kernels are
the \emph{primitive walls}.  A source mask \(M\) is \emph{sharp} when it
selects the maximum possible number of positive nonzero \(\rho\)-orbits; its
relative open realization cone is
\[
 K(M)=\{h\in V^-_\sigma:\ell_o(h)>0\ (o\in M),\quad
                  \ell_o(h)<0\ (o\notin M\cup Z)\},
\]
where \(Z\) is the structural-zero set.  A \emph{chamber} is the positive
projectivization \(K(M)/\mathbb R_{>0}\) of a nonempty sharp cone.

Two labelled chambers are \emph{raw sharp facet neighbours} when the
topological closures of their cones share a relative codimension-one face.
The certificate used here supplies exact interior potentials on both sides
whose joining segment meets exactly one primitive wall and no other wall.
The full centralizer \(G=C_{S_q}(\sigma)\) relabels masks and cones; its
vertex quotient has the chamber orbits \(GM\), while raw edges retain their
labelled endpoints and crossed wall.  At \(q=6\), \(G\cong C_2\wr S_3\) has
order 48.  A \emph{candidate flip} switches the complete sign class of one
of the 18 primitive walls having one form in each orientation.  It becomes a
realized neighbour only when the flipped sharp mask has a certified nonempty
cone and the exact one-wall segment test passes.

A bounded diagnostic does not enter the preceding proofs.  At \(m=3\), two
certified sharp facet adjacencies have endpoint defects
\((0,3)\) and \((4,1)\).  Thus defect is not 1-Lipschitz, and the maximum
facet jump is at least three.  A bounded panel of 90 candidate flips finds 23
realized neighbours, with jump histogram
\[
0:13,\qquad1:4,\qquad2:4,\qquad3:2.
\]
This panel is not the global facet graph and does not determine connectivity
or the maximum jump.  The obstruction of Section~\ref{sec:obstruction}
explains 24 of the 42 defect-four symmetry classes at \((3,0)\); it does not
classify the other 18 or provide an all-parameter defect formula.

\appendix
\section{Reproducibility and artifact boundary}
\label{app:reproducibility}

The finite seed certificates are reproduced and independently checked by
\begin{quote}\ttfamily\footnotesize\raggedright
python -B companion/verify\_five\_defect\_witnesses.py\par
python -B -O companion/verify\_five\_defect\_witnesses.py
\end{quote}
These five certificates prove attainment.  The reverse inclusion
\(\Dspec_{3,0}\subseteq\{0,1,2,3,4\}\) comes instead from the complete
six-letter verification described in Section~\ref{sec:seeds}.  Population
completeness is first reconstructed from the 36 local states and all 46,656
source triples by
\begin{quote}\ttfamily\footnotesize\raggedright
python -B companion/verify\_population.py\par
python -B -O companion/verify\_population.py
\end{quote}
The checker enumerates 987 source orbits and 1,759 canonical representative
masks, expands them to 83,736 labelled masks, and consults the tracked mask
rows only for a final exact-set comparison.  It imports neither Route A nor a
stored population as an enumeration oracle; it deliberately reuses the frozen
quotient-native Route B kernel rather than claiming a second implementation.

The portable
bundle contains opposed dynamic-programming tables for all 1,759
closure-centralizer classes.  The following standard-library checker
reconstructs one authenticated canonical closure per class, expands the 1,759
representative mask orbits under all 48 centralizer actions, verifies the
mask-to-class assignment of every one of the 83,736 tracked masks, checks all
3,518 tables against their Bellman recurrences, checks route agreement and
explicit optimal orders, and derives the displayed spectrum:
\begin{quote}\ttfamily\footnotesize\raggedright
python -B companion/verify\_bellman\_bundle.py\par
python -B -O companion/verify\_bellman\_bundle.py
\end{quote}
The public replay uses no historical SQLite database.

The four-letter rank-three obstruction of
Proposition~\ref{prop:rank-three-obstruction} is reproduced by
\begin{quote}\ttfamily\footnotesize\raggedright
python -B companion/verify\_rank\_three\_obstruction.py\par
python -B -O companion/verify\_rank\_three\_obstruction.py
\end{quote}
The standard-library checker enumerates all 99,000 linear extensions,
reconstructs the 176 attained bad-cell sets, verifies that forced cells and
pairwise incompatibilities are absent, and recovers the displayed triple as
the unique inclusion-minimal obstruction.

The exact occurrence statement in Proposition~\ref{prop:obstruction-occurrence}
is reproduced by
\begin{quote}\ttfamily\footnotesize\raggedright
python -B companion/verify\_ten\_relation\_core.py\par
python -B -O companion/verify\_ten\_relation\_core.py
\end{quote}
The checker reconstructs each canonical closure directly from the tracked
75-bit source masks and tests the core on all 1,759 work identities.  It
discovers 24 classes and all their
1,152 labelled members.  The human \(2+1+1=4\) lower bound and two directly
checked score-91 topological orders per discovered class give matching lower
and upper bounds.

The bounded facet diagnostic in Section~\ref{sec:diagnostics} is reproduced by
\begin{quote}\ttfamily\footnotesize\raggedright
python -B companion/verify\_facet\_counterexample.py\par
python -B -O companion/verify\_facet\_counterexample.py
\end{quote}
These commands certify the two displayed sharp adjacencies and the stated
90-flip panel only.  They do not construct the full six-letter facet graph.
The exact file manifest, prerequisites, expected outputs, PDF build, and
clean-clone procedure are given in \texttt{REPRODUCE.md} at repository root.

\bibliographystyle{plain}
\bibliography{references}

\end{document}
